\documentclass[a4paper,12pt,twoside]{memoir}

% Castellano
\usepackage[spanish,es-tabla]{babel}
\selectlanguage{spanish}
\usepackage[utf8]{inputenc}
\usepackage[T1]{fontenc}
\usepackage{lmodern} % scalable font
\usepackage{microtype}
\usepackage{placeins}

\RequirePackage{booktabs}
\RequirePackage[table]{xcolor}
\RequirePackage{xtab}
\RequirePackage{multirow}

% Links
\usepackage[colorlinks]{hyperref}
\hypersetup{
	allcolors = {red}
}

% Ecuaciones
\usepackage{amsmath}

% Rutas de fichero / paquete
\newcommand{\ruta}[1]{{\sffamily #1}}

% Párrafos
\nonzeroparskip


% Imagenes
\usepackage{graphicx}
\newcommand{\imagen}[2]{
	\begin{figure}[!h]
		\centering
		\includegraphics[width=0.9\textwidth]{#1}
		\caption{#2}\label{fig:#1}
	\end{figure}
	\FloatBarrier
}

\newcommand{\imagenflotante}[2]{
	\begin{figure}%[!h]
		\centering
		\includegraphics[width=0.9\textwidth]{#1}
		\caption{#2}\label{fig:#1}
	\end{figure}
}



% El comando \figura nos permite insertar figuras comodamente, y utilizando
% siempre el mismo formato. Los parametros son:
% 1 -> Porcentaje del ancho de página que ocupará la figura (de 0 a 1)
% 2 --> Fichero de la imagen
% 3 --> Texto a pie de imagen
% 4 --> Etiqueta (label) para referencias
% 5 --> Opciones que queramos pasarle al \includegraphics
% 6 --> Opciones de posicionamiento a pasarle a \begin{figure}
\newcommand{\figuraConPosicion}[6]{%
  \setlength{\anchoFloat}{#1\textwidth}%
  \addtolength{\anchoFloat}{-4\fboxsep}%
  \setlength{\anchoFigura}{\anchoFloat}%
  \begin{figure}[#6]
    \begin{center}%
      \Ovalbox{%
        \begin{minipage}{\anchoFloat}%
          \begin{center}%
            \includegraphics[width=\anchoFigura,#5]{#2}%
            \caption{#3}%
            \label{#4}%
          \end{center}%
        \end{minipage}
      }%
    \end{center}%
  \end{figure}%
}

%
% Comando para incluir imágenes en formato apaisado (sin marco).
\newcommand{\figuraApaisadaSinMarco}[5]{%
  \begin{figure}%
    \begin{center}%
    \includegraphics[angle=90,height=#1\textheight,#5]{#2}%
    \caption{#3}%
    \label{#4}%
    \end{center}%
  \end{figure}%
}
% Para las tablas
\newcommand{\otoprule}{\midrule [\heavyrulewidth]}
%
% Nuevo comando para tablas pequeñas (menos de una página).
\newcommand{\tablaSmall}[5]{%
 \begin{table}
  \begin{center}
   \rowcolors {2}{gray!35}{}
   \begin{tabular}{#2}
    \toprule
    #4
    \otoprule
    #5
    \bottomrule
   \end{tabular}
   \caption{#1}
   \label{tabla:#3}
  \end{center}
 \end{table}
}

%
%Para el float H de tablaSmallSinColores
\usepackage{float}

%
% Nuevo comando para tablas pequeñas (menos de una página).
\newcommand{\tablaSmallSinColores}[5]{%
 \begin{table}[H]
  \begin{center}
   \begin{tabular}{#2}
    \toprule
    #4
    \otoprule
    #5
    \bottomrule
   \end{tabular}
   \caption{#1}
   \label{tabla:#3}
  \end{center}
 \end{table}
}

\newcommand{\tablaApaisadaSmall}[5]{%
\begin{landscape}
  \begin{table}
   \begin{center}
    \rowcolors {2}{gray!35}{}
    \begin{tabular}{#2}
     \toprule
     #4
     \otoprule
     #5
     \bottomrule
    \end{tabular}
    \caption{#1}
    \label{tabla:#3}
   \end{center}
  \end{table}
\end{landscape}
}

%
% Nuevo comando para tablas grandes con cabecera y filas alternas coloreadas en gris.
\newcommand{\tabla}[6]{%
  \begin{center}
    \tablefirsthead{
      \toprule
      #5
      \otoprule
    }
    \tablehead{
      \multicolumn{#3}{l}{\small\sl continúa desde la página anterior}\\
      \toprule
      #5
      \otoprule
    }
    \tabletail{
      \hline
      \multicolumn{#3}{r}{\small\sl continúa en la página siguiente}\\
    }
    \tablelasttail{
      \hline
    }
    \bottomcaption{#1}
    \rowcolors {2}{gray!35}{}
    \begin{xtabular}{#2}
      #6
      \bottomrule
    \end{xtabular}
    \label{tabla:#4}
  \end{center}
}

%
% Nuevo comando para tablas grandes con cabecera.
\newcommand{\tablaSinColores}[6]{%
  \begin{center}
    \tablefirsthead{
      \toprule
      #5
      \otoprule
    }
    \tablehead{
      \multicolumn{#3}{l}{\small\sl continúa desde la página anterior}\\
      \toprule
      #5
      \otoprule
    }
    \tabletail{
      \hline
      \multicolumn{#3}{r}{\small\sl continúa en la página siguiente}\\
    }
    \tablelasttail{
      \hline
    }
    \bottomcaption{#1}
    \begin{xtabular}{#2}
      #6
      \bottomrule
    \end{xtabular}
    \label{tabla:#4}
  \end{center}
}

%
% Nuevo comando para tablas grandes sin cabecera.
\newcommand{\tablaSinCabecera}[5]{%
  \begin{center}
    \tablefirsthead{
      \toprule
    }
    \tablehead{
      \multicolumn{#3}{l}{\small\sl continúa desde la página anterior}\\
      \hline
    }
    \tabletail{
      \hline
      \multicolumn{#3}{r}{\small\sl continúa en la página siguiente}\\
    }
    \tablelasttail{
      \hline
    }
    \bottomcaption{#1}
  \begin{xtabular}{#2}
    #5
   \bottomrule
  \end{xtabular}
  \label{tabla:#4}
  \end{center}
}



\definecolor{cgoLight}{HTML}{EEEEEE}
\definecolor{cgoExtralight}{HTML}{FFFFFF}

%
% Nuevo comando para tablas grandes sin cabecera.
\newcommand{\tablaSinCabeceraConBandas}[5]{%
  \begin{center}
    \tablefirsthead{
      \toprule
    }
    \tablehead{
      \multicolumn{#3}{l}{\small\sl continúa desde la página anterior}\\
      \hline
    }
    \tabletail{
      \hline
      \multicolumn{#3}{r}{\small\sl continúa en la página siguiente}\\
    }
    \tablelasttail{
      \hline
    }
    \bottomcaption{#1}
    \rowcolors[]{1}{cgoExtralight}{cgoLight}

  \begin{xtabular}{#2}
    #5
   \bottomrule
  \end{xtabular}
  \label{tabla:#4}
  \end{center}
}




\graphicspath{ {./img/} }

% Capítulos
\chapterstyle{bianchi}
\newcommand{\capitulo}[2]{
	\setcounter{chapter}{#1}
	\setcounter{section}{0}
	\chapter*{#2}
	\addcontentsline{toc}{chapter}{#2}
	\markboth{#2}{#2}
}

% Apéndices
\renewcommand{\appendixname}{Apéndice}
\renewcommand*\cftappendixname{\appendixname}

\newcommand{\apendice}[1]{
	%\renewcommand{\thechapter}{A}
	\chapter{#1}
}

\renewcommand*\cftappendixname{\appendixname\ }

% Formato de portada
\makeatletter
\usepackage{xcolor}
\newcommand{\tutor}[1]{\def\@tutor{#1}}
\newcommand{\course}[1]{\def\@course{#1}}
\definecolor{cpardoBox}{HTML}{E6E6FF}
\def\maketitle{
  \null
  \thispagestyle{empty}
  % Cabecera ----------------
\noindent\includegraphics[width=\textwidth]{cabecera}\vspace{1cm}%
  \vfill
  % Título proyecto y escudo informática ----------------
  \colorbox{cpardoBox}{%
    \begin{minipage}{.8\textwidth}
      \vspace{.5cm}\Large
      \begin{center}
      \textbf{TFG del Grado en Ingeniería Informática}\vspace{.6cm}\\
      \textbf{\LARGE\@title{}}
      \end{center}
      \vspace{.2cm}
    \end{minipage}

  }%
  \hfill\begin{minipage}{.20\textwidth}
    \includegraphics[width=\textwidth]{escudoInfor}
  \end{minipage}
  \vfill
  % Datos de alumno, curso y tutores ------------------
  \begin{center}%
  {%
    \noindent\LARGE
    Presentado por \@author{}\\ 
    en Universidad de Burgos --- \@date{}\\
    Tutor: \@tutor{}\\
  }%
  \end{center}%
  \null
  \cleardoublepage
  }
\makeatother


% Datos de portada
\title{título del TFG \\Documentación Técnica}
\author{nombre alumno}
\tutor{nombre tutor}
\date{\today}

\begin{document}

\maketitle



\cleardoublepage



%%%%%%%%%%%%%%%%%%%%%%%%%%%%%%%%%%%%%%%%%%%%%%%%%%%%%%%%%%%%%%%%%%%%%%%%%%%%%%%%%%%%%%%%



\frontmatter


\clearpage

% Indices
\tableofcontents

\clearpage

\listoffigures

\clearpage

\listoftables

\clearpage

\mainmatter

\appendix

\apendice{Plan de Proyecto Software}

\section{Introducción}

\section{Planificación temporal}

\section{Estudio de viabilidad}

\subsection{Viabilidad económica}

\subsection{Viabilidad legal}



\apendice{Especificación de Requisitos}

\section{Introducción}

\section{Objetivos generales}

\section{Catalogo de requisitos}

\section{Especificación de requisitos}



\apendice{Especificación de diseño}

\section{Introducción}
A continuación especificaremos la forma en que hemos organizado los elementos del proyecto, así como el porqué.

\section{Diseño de datos}
\subsection{Puntos}
Necesitamos una convención para que nuestro sistema de importar y exportar puntos funcione correctamente. El archivo \textit{JSON} (ver código \ref{JSON-schema}) se compone por:
\begin{itemize}
	\item \texttt{filename} El nombre del modelo al que pertenece.
	\item \texttt{annotations} Una lista con las diferentes anotaciones.
	\item \texttt{measurements} Una lista con las medidas.
	\item \texttt{timestamp} Una marca de tiempo de cuando se ha validado en el servidor.
	\item \texttt{checksum} Un código de comprobación (\textit{md5} para ser concretos.)
\end{itemize}
Cada una de las anotaciones se compone por una etiqueta y un punto. Las medidas son una etiqueta y dos puntos. Finalmente, los puntos son tres números de tipo \textit{float} (\texttt{x}, \texttt{y}, \texttt{z}).

La estructura de los archivos \textit{JSON} puede verse en el código \ref{JSON-schema}.

\begin{lstlisting}[language=json, float, caption=Esquema JSON, label=JSON-schema]
{"filename": "a filename",
 "annotations": [annotation, ..., annotation],
 "measurements": [measurement, ..., measurement]}

annotation = {"tag": "a tag",
              points: [point]}
		
measurement = {"tag": "a tag",
               points: [point, point]}
		  
point = {"x": float,
         "y": float,
         "z": float}
\end{lstlisting}


\section{Diseño procedimental}
En el momento que el servidor está en línea, éste puede aceptar peticiones de inicio de sesión. Cuando un usuario lo solicita, es redirigido a la página de \textit{login}, en el que tendrá que introducir su correo y contraseña de UBUVirtual. Tras ello, el servidor buscará en su base de datos de usuarios si éste está presente en ella. Si no está, redirigiremos al usuarios a la página de \textit{login} anteriormente mencionada. Si pertenece, entonces el servidor lanza una consulta a la \textit{API} de UBUVirtual, para saber si también se encuentra reconocido como usuario en dicha plataforma. Si no consta el correo, o si la contraseña es incorrecta, se redirige a la página de \textit{login}. En caso de que ambas preguntas sean correctas, finalmente comprobaremos si había entrado con la intención de abrir un modelo concreto en caso afirmativo, le redirigiremos a aquel que hubiera solicitado; si no, visualizará la estantería de modelos.

En la figura \ref{fig:flux-diagram-login} se aprecia cómo sucede el proceso de \textit{login}.
\imagenpersonalizada{flux-diagram-login}{Proceso de \textit{login}}{0.75}

Por orden cronológico, el siguiente proceso que podremos realizar será elegir el modelo, tal y como se aprecia en la figura \ref{fig:flux-diagram-model-selection}.
\imagenpersonalizada{flux-diagram-model-selection}{Selección de modelo}{0.75}

\section{Diseño arquitectónico}
En la figura \ref{fig:class-diagram} el diagrama de clases que define la parte JavaScript del proyecto. La clase \texttt{<<Utils>>} aunque separada del resto, sí está relacionada con las demás clases, aunque para facilitar la comprensión hemos evitado las uniones. Dependen de ella las clases \texttt{<<AnnotationTool>>}, \texttt{<<MeasurementTool>>}, \texttt{<<PointManager>>} y \texttt{<<Scene>>}.
\imagenpersonalizada{class-diagram}{Diagrama de clases de parte JavaScript}{1.0}

Sin embargo, no añadiremos el diagrama de paquetes JavaScript, puesto que solamente existe uno.



\apendice{Documentación técnica de programación}

\section{Introducción}
A continuación documentaremos la aplicación tanto para alguien que vaya a realizar su despliegue~\ref{manual-despliegue}, como para aquel que la vaya a modificar (ver sección~\ref{manual-programador}).

\section{Manual de despliegue}\label{manual-despliegue}
Vamos a ver los pasos a realizar para desplegar la aplicación desde cero. No obstante y con el fin de facilitar el despliegue, tenemos un par de \textit{scripts} para dicha tarea: 
podemos emplear \texttt{<<WebPage/install\_run.bat>>} para Windows, además de \texttt{<<WebPage/install\_run.sh>>} para Linux. Si la ejecución del \textit{script} lanza algún error, siempre podemos realizar una instalación paso a paso.

\subsection{Lanzar un terminal o consola}
Para comenzar, lo primero que tendremos que hacer es ejecutar una consola, algo realmente fácil. En Windows por ejemplo, pulsando \texttt{<<Windows + R>>} podemos ejecutar un comando. En este caso, queremos lanzar la consola de comandos, así que teclearemos \texttt{<<cmd>>} en el recuadro de texto y pulsaremos ejecutar. Así de fácil.
En versiones recientes de Windows, también podremos utilizar la \texttt{<<Powershell>>} si así lo deseamos. En el menú de inicio podemos buscar la aplicación como tal.
En Linux lanzar el terminal incluso suele estar asociado a un atajo de teclado: \texttt{<<Ctrl + Alt + T>>}. En caso contrario, podemos buscar en el menú de aplicaciones, suele estar nombrado como \texttt{<<Terminal>>}, otras veces como \texttt{<<XTerm>>}, etc.
Finalmente, en Mac se llama \texttt{<<Terminal>>}.

\subsection{Instalando Python}\label{instalar-python}
Los componentes en que se basa el servidor web necesitan Python, que será nuestra primera parada. Para ello, iremos a la página de descargas de Python, y nos descargaremos la versión \texttt{<<3.5.2>>}\footnote{\url{https://www.python.org/downloads/release/python-352/}} para nuestro sistema operativo. Tras la descarga, procederemos a su instalación. Durante el proceso en Windows, se preguntará si desea añadir Python al \texttt{PATH} del sistema. Si desea utilizar los comandos tal cual se escriben aquí, debe aceptar esta opción. Si no, en vez de escribir \texttt{<<python3>>} o \texttt{<<pip3>>}, deberá introducir la ruta absoluta de los binarios.

\subsection{Creando un entorno virtual}\label{crear-venv}
Tras completar el paso anterior, vamos a instalar algunos componentes. Si ya tenemos más instalaciones de Python, o si simplemente queremos que nuestra instalación siga limpia, deberemos hacer un pequeño rodeo creando un  entorno virtual sobre el que recaigan los componentes. Insistimos, si has leído estos casos y no te sientes identificado, puedes saltarse esta explicación.

\noindent\verb|$ python3 -m venv mi_venv|

Lo que hace el código arriba descrito\footnote{Es posible que el comando \texttt{<<python3>>} no sea reconocido en ocasiones; sería posible que el alias real para la instalación sea \texttt{<<python>>.}} es crear un entorno virtual en la carpeta actual bajo la carpeta visor. Una vez creado, queremos que las dependencias se instalen en él así que necesitamos activarlo primero. Para ello, escogemos la opción según nuestro sistema operativo, escogiéndola de la tabla~\ref{tabla:comandos-source} y nuestra consola para ejecutarla en un terminal.

\begin{tabular}{lll}
		\otoprule
	Posix   & bash/zsh   & \verb|$ source mi_venv/bin/activate|          \\
		\cline{2-3}
	& fish       & \verb|$ . mi_venv/bin/activate.fish|          \\
		\cline{2-3}
	& csh/tcshq  & \verb|$ source mi_venv/bin/activate.csh|      \\
	\midrule
	Windows & cmd.exe    & \verb|C:\> mi_venv\Scripts\activate.bat|    \\
		\cline{2-3}
	& Powershell & \verb|PS C:\> mi_venv\Scripts\Activate.ps1| \\
		\bottomrule
\end{tabular}\label{tabla:comandos-source}

Ahora, si instalamos las dependencias se hará dentro de este entorno, sin <<manchar>> nuestra instalación original. Cuando necesitemos salir de dicho entorno virtual, teclearemos \texttt{<<deactivate>>}, o simplemente cerraremos el terminal.



\subsection{Instalando las dependencias Python}\label{instalar-dependencias-python}
En la carpeta \texttt{<<Webpage>>}, encontramos un archivo \texttt{<<requirements.txt>>}. A éste debemos hacer mención con el comando\footnote{Con el comando \texttt{<<pip3>>} es posible que también necesitemos escribirlo sin el \texttt{<<3>>}.}:

\noindent\verb|$ pip3 install -r <ruta proyecto>/Webpage/requirements.txt|

Con las dependencias ya instaladas, tenemos que dar permisos de acceso a los usuarios que deseemos.

\subsection{Dando permisos a usuarios}
Si queremos que un usuario pueda utilizar la aplicación, tenemos que incluirlo en la base de datos. Con el fin de facilitar añadir usuarios a dicha tabla, tenemos un método del que nos podemos aprovechar en el módulo \texttt{<<MyApp>>}, llamado \texttt{<<import\_users\_to\_db>>}. Para llamarlo, desde la línea de comandos debemos seleccionar la carpeta \texttt{<<WebPage>>} mediante el comando \texttt{<<cd>>}. Como argumento al mismo, escribiremos la carpeta hacia el cual nos queremos dirigir. Si lo que deseamos es ir un directorio hacia arriba, escribiremos como parámetro \texttt{<<..>>}.

Si la consola nos indica por ejemplo que estamos sobre la carpeta MyApp del proyecto, subiremos una posición escribiendo:

\noindent\verb|$ cd ..|

Si la consola nos indica que estamos en la carpeta que contiene el proyecto (típicamente \texttt{<<3D-Viewer>>)}, entonces bajaríamos mediante:

\noindent\verb|$ cd WebPage|

Nótese que podríamos encadenar varios de ellos. Por ejemplo, si estuviésemos en la carpeta `Documentación', entonces:

\noindent\verb|$ cd ../WebPage|

Volviendo a la explicación, lo que tenemos que hacer es cambiar la carpeta seleccionado en la consola.

\noindent\verb|$ cd <ruta proyecto>/WebPage|

Una vez hecho esto, podemos inicial el terminal interactivo de python mediante:

\noindent\verb|$ python3|

Dentro del mismo, podremos importar el módulo \texttt{<<MyApp>>} mediante:

\noindent\verb|>>>import MyApp|

Ahora, emplearemos el método que queríamos, \texttt{<<import\_users\_to\_db>>}. Podemos hacerlo sin parámetro:

\noindent\verb|>>>MyApp.import_users_to_db()|

Y, entonces buscará el archivo \texttt{<<WebPage\slash instance\slash users\textunderscore to\textunderscore import.csv>>} por defecto. Si por el contrario, deseamos que busque otro archivo (siempre dentro de la carpeta \texttt{<<instance>>}), se lo podemos pasar como parámetro:

\noindent\verb|>>>MyApp.import_users_to_db('otro_archivo')|

De esta manera, los usuarios dentro de este archivo serán importados a la base de datos para otorgarles permiso de acceso. El archivo que vamos a emplear para decirle los usuarios, tiene que tener un formato especial: \textit{CSV}. Éste es realmente sencillo de realizar mediante un programa de hoja de cálculo como <<Excel>>, <<Numbers>> o <<Calc>>, por mencionar los más conocidos. Simplemente llenaremos las celdas de la siguiente manera:

\tablaSmallSinColores{Ejemplo de \textit{CSV}}{ccc}{ejemplo-csv}
{
	\textbf{email} & \textbf{nombre} & \textbf{rol} \\
}
{
	email@mail.com & pepito & POTUS \\
	secret@mail.com & 007 Agent & slave \\
}

Como se ve, en cada columna hay un tipo de dato, y el nombre de dicho dato se determina mediante la cabecera de dicha columna. Las celdas inferiores a la misma contendrán los datos de dicho valor para cada uno de los usuarios. Lo mismo pasa con el resto de las columnas.

Pista: para introducir estos datos en la hoja de cálculo correspondiente, normalmente se pueden pegar directamente desde una tabla en un navegador, etc.

Una vez los datos en la hoja de cálculo el programa correspondiente tendrá un pequeño asistente para conseguir exportar a \textit{CSV} los valores introducidos. Dicha orden suele encontrarse en el menú \texttt{<<Archivo>>}, bajo la opción \texttt{<<Exportar...>>}, \texttt{<<Exportar como...>>}, \texttt{<<Guardar como...>>}, etc. Normalmente podremos emplear los parámetros por defecto del asistente, pero no nos olvidemos de guardar el archivo en la carpeta \texttt{<<instance>>} con el nombre adecuado.

Nota: En el momento de redactar esta sección, el único campo necesario es el \textit{email}, pero para saber a ciencia cierta cuáles son los campos por almacenar, conviene echar un vistazo al \textit{script} de creación en \texttt{<<WebPage\slash MyApp/sql\textunderscore scripts\slash schema.sql>>}.

\subsection{Lanzando el servidor}
Ahora, bajo la carpeta \texttt{<<Webpage>>} encontraremos el fichero \texttt{<<runserver.py>>} al que haremos referencia con la orden:

\noindent\verb|$ python3 <ruta proyecto>/Webpage/runserver.py|

Et \textit{voilà}: el servidor está lanzado, y el terminal nos indica la dirección que debemos introducir para consultar el contenido: \texttt{<<localhost:5000>>} o \texttt{<<127.0.0.1:5000>>} como se prefiera. Por fin debería ver la interfaz de la aplicación.



\section{Manual del programador}\label{manual-programador}
Veamos los diferentes elementos que nos pueden ser de utilidad a la hora de realizar cambios en el proyecto.

\subsection{Instalando dependencias Node}
Para labores de desarrollo, necesitaremos tener instalado Node y algunas dependencias extra. Podremos si tecleamos desde la carpeta \texttt{<<TFG/WebPage/>>} el siguiente comando:

\noindent\verb|$ npm install|

Éste y todos los comandos que sean con npm, deberán ser ejecutados desde la misma ruta.

\subsection{Transpilación JavaScript}
Durante el desarrollo del proyecto, separamos el código JavaScript en diferentes fuentes, que requieren para funcionar de realizar una transpilación a un estándar que comprendan los navegadores actuales. Aunque éstos soportan gran parte de las características de \textit{ES6} (\textit{ES2015}), estándar en el que hemos escrito el fuente original, hay elementos que todavía no soportan, como los \textit{import} entre las diferentes clases. Adicionalmente, aprovechamos para hacer el archivo destino más pequeño, en un proceso denominado \textit{minifying}.
Para ello, vamos a empaquetar el código de una manera más adecuada, y nos ayudaremos de rollup:

\noindent\verb|$ npm run build|

Si por ejemplo estamos desarrollando, sería molesto llamar a esta función cada vez que hiciésemos un cambio. Para ello, hemos implementado una configuración que nos permitirá invocar una sola vez a un comando de transpilación, y cada vez que realicemos un cambio automáticamente se generará el fichero destino en formato \textit{minified}. Para hacer esto, tenemos que llamar al comando:

\noindent\verb|$ npm run build dev|

\subsection{Construir la documentación JavaScript}
Para hacerlo, tenemos otro \textit{script} incluido en la configuración de Node para facilitarnos la tarea. Solamente hay que escribir:

\noindent\verb|$ npm run jsdoc|

\subsection{Construir la documentación de Python}
Construir la documentación de Python es igual de sencillo:

\noindent\verb|$ npm run pydoc-http|

Eso sí, hay que tener en cuenta que éste comando no nos avisa de que sigamos algunas pautas recomendables, así que necesitamos de otro comando para revisar la sintaxis:

\noindent\verb|$ npm run pydoc-style|

\subsection{Usando el analizador de modelos: eliminando puntos superfluos}
Con el fin de mejorar los modelos, añadimos varias funciones a nuestras herramientas con el fin de <<recortar>> en la medida de lo posible el espacio que ocupan los modelos, además de evitar errores posteriores con los mismos.

Hemos empaquetado toda la funcionalidad en un módulo Python, por lo que abriremos la dicha consola sobre la carpeta \texttt{<<Resources>>}:

\noindent\verb|$ python3|

A continuación, cargaremos la función a emplear:

\noindent\verb|>>> from ply_utils import process_ply_file|

Con la función cargada, aplicaremos los parámetros necesarios.

\noindent\verb|>>> process_ply_file('una ruta de archivo', 'encoding')|

La ruta de archivo puede ser tanto relativa como absoluta.
El parámetro de \texttt{encoding} es opcional, y si se incluye debe ser \texttt{<<b'ascii'>>}, \texttt{<<b'binary\_little\_endian'>>} o \texttt{<<b{'binary\_big\_endian'}>>}.

Este último parámetro es el que nos permitirá comprimir los archivos \textit{ASCII} en formato binario, más compacto. Así mismo, podremos generar un fichero \textit{ASCII} a partir de uno binario.

\subsection{Estructura de directorios}
\noindent Para comenzar, mostraremos los archivos existentes en la raíz del proyecto:
\begin{itemize}
	\item \textbf{.codeclimate.yml} Contiene la configuración para calidad del código fuente de CodeClimate\footnote{\url{https://codeclimate.com/github/Alberto-Vivar/3D-Viewer}}, un comprobador automático de código.
	\item \textbf{.gitignore} Posee los diferentes elementos que el sistema \textit{GIT} debe evitar revisar.
	\item \textbf{README.md} Archivo léeme en formato \textit{Markdown} que aparece en la portada de GitHub.
\end{itemize}

Ahora mostraremos las carpetas disponibles en la raíz del proyecto:
\begin{itemize}
	\item Documentación
	\item WebApp
	\item Resources
\end{itemize}

\noindent A continuación las describiremos en mayor detalle.

\subsubsection{Documentación}
En ella tendremos el conjunto de documentaciones generadas para el proyecto, cuyas carpetas son:
\begin{itemize}
	\item \textbf{javascript} En ella tenemos la documentación generada para la parte Javascript, esto es, toda la lógica e interfaz del visor.
	\item \textbf{LaTex} Contiene la memoria y anexos del proyecto.
	\item \textbf{python} Posee la documentación generada mediante Sphinx\footnote{\url{http://www.sphinx-doc.org/en/stable/}} para el servidor.
\end{itemize}

\subsubsection{WebApp}
En ella se encuentra definido tanto el servidor como la página web en si misma. Veamos los directorios que la componen:
\begin{itemize}
	\item \textbf{instance} Aunque no es un directorio incluido en el repositorio, debemos mencionarlo, puesto que en él incluiremos algunas configuraciones propias de la instancia que estemos desplegando o desarrollando y que no queramos incluir en la versión que subamos al servidor. Algunos archivos de utilidad que meteremos son:
		\begin{itemize}
			\item \textbf{config.py} En este archivo incluiremos algunos parámetros de configuración que debemos sobrescribir, como el estado de \texttt{<<DEBUG>>}, o el secreto de aplicación \texttt{<<SECRET\_KEY>>}.
			\item \textbf{users\_to\_import.csv} Es la ruta por defecto para el archivo \textit{CSV} que nos permite importar la lista de usuarios a los que queremos dar permiso.
		\end{itemize}
	\item \textbf{MyApp} En él está la aplicación en sí misma, casi toda la lógica de negocio y toda la visualización.
		\begin{itemize}
			\item \textbf{sql\_scripts} Incluye los \textit{scripts} de utilidad para la base de datos.
			\item \textbf{static} Directorio para cargar elementos estáticos de las páginas.
			\begin{itemize}
				\item \textbf{js} \textit{Scripts} que importaremos. Aquí tendremos el código JavaScript una vez transpilado, y por tanto, gran parte de la lógica del visor.
				\item \textbf{uploads} En este directorio terminan los modelos, además de los \textit{PNG} como sus miniaturas.
			\end{itemize}
			\item \textbf{templates} En él se encuentran las plantillas en formato jinja2\footnote{\url{http://jinja.pocoo.org/}}.
			\item \textbf{translations} En este directorio encontraremos los ficheros plantilla para poder realizar la traducción con Babel.
			\item \textbf{babel.cfg} Configuración con la que Babel busca las etiquetas identificadoras.
			\item \textbf{forms.py} Contiene constructores de los diferentes formularios.
			\item \textbf{\_\_init\_\_.py} Constructor del módulo.
			\item \textbf{messages.pot} Binario con las traducciones de Babel ya realizadas
			\item \textbf{User.py} Clase necesaria para la gestión de usuarios.
			\item \textbf{users.db} Base de datos que alberga la información de los usuarios admitidos.
			\item \textbf{views.py} Contiene los diferentes \textit{endpoints} de nuestro servidor.
		\end{itemize}
	\item \textbf{src} Código fuente de JavaScript sin compilar.
	\item \textbf{typings} Archivos de configuración para nuestro \textit{IDE} (Visual Studio Code).
\end{itemize}

Ahora, veamos los archivos que contiene:
\begin{itemize}
	\item \textbf{config.py} Archivo de configuración predeterminada. Estos parámetros se verán sobrescritos por los añadidos en \texttt{instance/config.py}.
	\item \textbf{install\_run.bat} \textit{Script} para Windows que genera un entorno virtual, lo registra, instala las dependencias Python, y lanza el servidor.
	\item \textbf{install\_run.sh} El mismo \textit{script}, pero para Linux.
	\item \textbf{jsconfig.json} Archivo de configuración para \textit{VSCode}.
	\item \textbf{package.json} Define las dependencias y comandos disponibles llamándolo desde npm (Node).
	\item \textbf{requirements.txt} Requisitos de los módulos para Python.
	\item \textbf{rollup.config.js} Configuración para el \textit{plugin} que empaqueta el código.
	\item \textbf{runserver.py} \textit{Script} que carga el módulo principal para lanzar el servidor.
	\item \textbf{typings.json} Configuración IntelliSense de Visual Studio Code.
\end{itemize}

\subsubsection{Resources}
Contiene algunos recursos que no sabíamos donde meter, pues no pertenecen estrictamente a ninguna de las partes anteriores. Entre sus elementos se encuentran los archivos <<fuente>> de los diagramas, o los generadores de informes para los archivos \textit{PLY}, que posiblemente integremos en un futuro sobre el servidor.

\subsection{Babel - Internacionalización}
Como el proceso ya está iniciado, en el manual detallaremos la actualización y no el comienzo.
La instalación de Babel es un requisito incluido en \texttt{<<requirements.txt>>}, así que podemos instalarlo junto con el resto de dependencias del proyecto (preferentemente en el entorno virtual):

\noindent\verb|$ pip3 install -r requirements.txt|

Los ficheros que Babel procesa están determinados por babel.cfg, así que cualquier necesidad extra deberemos reflejarla en dicho archivo.
Al añadir o eliminar las cadenas a procesar del código, necesitamos que Babel sepa cuales son sus cambios, así que le pedimos las extraiga:

\noindent\verb|$ pybabel extract -F babel.cfg -k lazy_gettext -o messages.pot .|

Nota: Ojo al punto final, no es una errata, tiene significado (directorio actual).

Aunque es improbable que dejemos de emplear las sustituciones perezosas (son las lazy\_gettext), si dejasen de ser necesarias podríamos quitar la parte \texttt{<<-k lazy\_gettext>>} del comando.
Si queremos añadir un idioma que todavía no existe, entonces tenemos que crear un fichero inicial del mismo:

\noindent\verb|$ pybabel init -i messages.pot -d translations -l es|

El parámetro \texttt{<<-d translations>>} dictamina que el nuevo fichero tendrá como destino la carpeta translations, para que tengamos todas las traducciones juntas.
El parámetro \texttt{<<-l es>>} determina el lenguaje; en este caso es el castellano, así que si queremos otro diferente tenemos que buscar la abreviatura adecuada\footnote{\url{http://babel.pocoo.org/en/latest/api/languages.html\#module-babel.languages}} según el proyecto que hemos empleado (flask-babel)\footnote{\url{https://github.com/python-babel/flask-babel}}\footnote{\url{https://pythonhosted.org/Flask-Babel/}}, tendríamos que revisar la versión de \textit{CLDR} empleada por Babel.

Si el idioma ya estaba creado, entonces no queremos que nos sobre-escriba lo que ya teníamos traducido. Entonces, le pediremos nos actualice las referencias mediante:

\noindent\verb|$ pybabel update -i messages.pot -d translations|

Añadirá las nuevas cadenas, y aquellas que ya no encuentre, cambiará su marca de \texttt{<<\#:>>} a \texttt{<<\#\textasciitilde>>}.

El único paso que nos quedará entonces será el de generar los ficheros que emplea Babel:

\noindent\verb|$ pybabel compile -d translations|

Y eso es todo, ya podemos tener las últimas etiquetas traducidas de forma dinámica.



\section{Problemas puntuales}\label{manual-programador-problemas-puntuales}
El error \texttt{<<No enviroment provided>>} salta a la hora de llamar activate si el directorio de trabajo actual de la consola es la carpeta bin del entorno virtual.
\newline
Solución: cambiar el directorio de trabajo al menos una carpeta más <<arriba>>.

Al crear el entorno virtual no encuentra pip.
\newline
Solución 1: usar la opción \texttt{<<\text{--without-pip}>>} durante la creación y posteriormente instalar \textit{PIP} en el entorno virtual.
\newline
\noindent\verb|$ sudo apt-get install python3-pip|
\newline
Solución 2: usar otra distribución de Python.

A la hora de generar los archivos necesarios para Babel, salta el siguiente error:
\newline
{\scriptsize\noindent\verb|UndefinedEnvironmentName: 'extra' does not exist in evaluation environment.|}
\newline
Solución: actualizar setuptools\footnote{\url{http://stackoverflow.com/questions/43556482/pybabel-extract-undefinedenvironmentname-extra-does-not-exist-in-evaluation-e}}



\section{Pruebas del sistema}
Por problemas de tiempo, y por no haber realizado el proyecto mediante la filosofía \textit{TDD}~\cite{wiki:TDD}, no se han realizado pruebas formales del sistema.

\apendice{Documentación de usuario}

\section{Introducción}
En esta sección vamos a explicar los diferentes aspectos involucrados en el uso de la aplicación.

\section{Requisitos de usuarios}
La parte del visor funciona en Firefox 50 y posteriores, Safari 10.1 y superior, y Chrome en versiones superiores a la 58.
Se requiere de conexión a Internet, además de estar autorizado a usar la aplicación.

\section{Instalación}
Suponemos que el usuario final tendrá un navegador web disponible, así que no necesitará instalar nada. No obstante, garantizamos la funcionalidad con Firefox\footnote{\url{https://support.mozilla.org/es/products/firefox/install-and-update-firefox}}.

\section{Manual del usuario}
A continuación explicaremos paso a paso cómo utilizar las funciones disponibles en la aplicación.

También definiremos algunos términos para conseguir una mejor comprensión de las instrucciones:
\begin{itemize}
	\item Barra de navegación: aparece en la parte superior de casi todas las vistas, y nos sirve para navegar por los diferentes elementos de la interfaz. Nos permite acceder al almacén de modelos, subir otros nuevos o cerrar sesión.
	\item Modelo: versión digital de cada elemento físico que vamos a emplear (restos óseos, fósiles, etc).
	\item Vista: cada una de las interfaces que es fundamentalmente diferente, en nuestro caso tenemos una página de \textit{login}, otra para subir nuevos modelos, otra más para previsualizar aquellos disponibles y otra para subir otros nuevos.
	\item Anotación: elemento gráfico que permite etiquetar un punto o zona del modelo mediante una esfera pegada al mismo. La etiqueta aparece en el menú de anotaciones sobre la barra de herramientas.
	\item Medida: elemento gráfico que permite conocer la distancia lineal entre dos puntos del modelos, además de poner un nombre concreto a dicha medida. En el modelo, se encuentra representada a través de dos esferas y una línea que las une. Aparecen sobre la barra de herramientas (concretamente sobre el menú de medidas) su etiqueta y medida.
	\item Barra de herramientas: aparece en la parte izquierda en la vista del visor, y permite realizar acciones con anotaciones y medidas. Nos da acceso a las diferentes herramientas: anotaciones, medidas, y asistente de importación\slash exportación.
	\item Visor: elemento gráfico que se encarga de mostrar un modelo. Es el rectángulo negro sobre el que se ve el modelo.
\end{itemize}

\subsection{Login del usuario}
Una vez estemos autorizados a usar la aplicación (típicamente el responsable de la asignatura nos lo permitirá), lo primero que necesitamos hacer es conectarnos a la misma. Para ello, abriremos un navegador \textit{web} compatible e introduciremos la dirección apropiada. Una vez allí, nos aparecerá la ventana de \textit{login}, similar a la figura ~\ref{fig:viewer-login}. En ella vamos a introducir las mismas credenciales de acceso que empleamos para los servicios \textit{on-line} en la Universidad de Burgos. Dicha página nos redirigirá automáticamente a la sección adecuada.
\imagenpersonalizada{viewer-login}{\textit{Login} de la aplicación}{1}

\subsection{Barra de navegación}
Excepto en la página de \textit{login}, tenemos un elemento común que aparece en todas las vistas de la aplicación. Podemos aprovecharnos de ella para realizar tres acciones: ir a la vista de modelos, ir a la vista para subir éstos, y cerrar la sesión actual. Podemos apreciar dicha barra en la figura~\ref{fig:viewer-navigation-bar}.
\imagenpersonalizada{viewer-navigation-bar}{Barra de navegación}{1}

\subsection{Repositorio de modelos}
Una vez se haya conectado, lo más probable es que la primera página que vea sea la estantería de modelos. En ella podrá elegir el modelo que desea visualizar; para ello, solamente debe pinchar sobre este y automáticamente se abrirá un visor con el modelo seleccionado. Los modelos aparecerán con una miniatura para ser fácilmente identificables, tal y como se aprecia en la figura~\ref{fig:viewer-shelf}.
\imagenpersonalizada{viewer-shelf}{Repositorio de modelos}{1}

\subsection{Manipulando el modelo}
Existen diferentes acciones básicas que podemos realizar para manipular el modelo y adecuar su posición consiguiendo así una mejor perspectiva del mismo; estas acciones son rotar, mover y enfocar.

El enfoque o \textit{zoom} es realmente sencillo, solamente necesitamos mover la rueda del ratón. Para que esto suceda, el puntero debe estar ubicado dentro del visor. Como acción alternativa a esta (si nuestra rueda no funciona tan bien como quisiéramos), también podemos pinchar y mantener pulsado mediante el botón central del ratón (la rueda) mientras movemos el ratón.

La rotación requiere que pinchemos con el botón izquierdo del ratón sobre el visor, mantengamos pulsado y movamos el puntero del ratón. Cuando soltemos, el movimiento del ratón dejará de afectar al modelo.

El desplazamiento es similar a la rotación, pero se diferencia en que usa el botón derecho del ratón en vez de el izquierdo.

\subsection{Acciones avanzadas con el modelo}
Solo la visualización del modelo nos da interesantes oportunidades para conocerlo, pero tenemos disponibles un mayor número de opciones.

Viendo la figura~\ref{fig:viewer-interface} podemos observar la disposición de la interfaz.
\imagenpersonalizada{viewer-interface}{Interfaz de usuario del visor}{1}

\subsubsection{Acciones comunes}
Tanto para el menú de anotaciones como el de medidas, existen acciones comunes. Éstas son:
\begin{itemize}
	\item Añadir: este botón añade un nuevo elemento tanto en el menú adecuado como en el visor.
	\item Editar: mediante dicho botón podemos cambiar la etiqueta del mismo.
	\item Eliminar: pulsando este botón borra la presencia de dicho elemento tanto del menú correspondiente como en el visor.
	\item Seleccionar/Deseleccionar: esta acción no se lleva a cabo pulsando ningún botón; se realiza haciendo \textit{click} sobre cualquier elemento (anotación o medida) que aparezca tanto en el visor como en la lista del menú correspondiente. Sobre las medidas, se realiza pulsando sobre cualquiera de los dos puntos que la componen.
	\item Deseleccionar todo: dicho botón nos permite desmarcar todos los elementos del mismo tipo que el del menú empleado.
	\item Cambiar la escala: dentro del visor, en la parte superior derecha, tenemos un pequeño menú que nos permitirá realizar diversos ajustes con la escala del modelo, además de visualizar una malla para hacernos una idea de las dimensiones.
\end{itemize}

\subsubsection{Anotaciones}
En la figura~\ref{fig:viewer-femur-annotation-greater-trocanter} podemos apreciar dos anotaciones sobre un fémur: el <<trocánter mayor>>, seleccionado y que aparece resaltado en la lista y en verde el modelo, y el <<trocánter menor>>, que no se encuentra seleccionado.
\imagenpersonalizada{viewer-femur-annotation-greater-trocanter}{Dos anotaciones sobre un fémur}{1}

Para añadir una nueva anotación, pincharemos primero en el botón <<Añadir>> del menú de anotaciones, y después pincharemos sobre el modelo. Actualmente no se puede cancelar la acción de añadir, así que si pinchamos accidentalmente dicho botón, y no queríamos crear ninguna anotación, primero tendremos que añadir un punto cualquiera, y después eliminarlo.

Para editar una anotación, primero tenemos que seleccionar la deseada. Seleccionaremos solamente una, ya sea desde el menú o a través del modelo. Después pulsaremos sobre <<Editar>>, y nos aparecerá una ventana para cambiar la etiqueta. Una vez cambiado el texto, podremos guardar los cambios pulsando <<Save>> o descartarlos, pudiendo pinchar tanto en <<Cancel>> como cerrando la ventana.
En la figura~\ref{fig:viewer-femur-annotation-editing} podemos observar el diálogo que nos permite editar la anotación.
\imagenpersonalizada{viewer-femur-annotation-editing}{Editando una anotación}{1}

Si lo que queremos es eliminar una anotación, seleccionaremos aquellas que deseemos descartar (en este caso sí pueden ser varias). Acto seguido, pulsaremos eliminar. Si no tenemos ninguna seleccionada antes de pulsar en <<Eliminar>>, podremos borrar la siguiente anotación sobre la que pulsemos. Desafortunadamente, en este segundo caso tenemos un comportamiento similar al de añadir: hasta que no eliminemos una anotación, no podemos cancelar este estado. Una <<solución>>, si nos hemos equivocado y no queremos eliminar ninguna de las anotaciones disponibles, reside en añadir nosotros otra de forma aleatoria. Así, la creamos con el procedimiento anterior, y después de ésto, lo que haremos es pinchar otra vez sobre la misma en el visor.

Finalmente, si tenemos varias anotaciones seleccionadas y queremos que no lo estén, únicamente pincharemos en el botón <<Deseleccionar todo>> en el menú correspondiente.

\subsubsection{Medidas}
Para las medidas, los procedimientos antes descritos son similares, así que explicaremos las diferencias a continuación (además de emplear los botones de su menú).

Cuando añadimos, tenemos que pulsar sobre dos ubicaciones del modelo para conseguir finalizar la medida.

Editar se hace del mismo modo.

Para eliminar una medida, la selección se puede realizar pinchando cualquiera de las dos esferas que componen una medida.

Y por supuesto, deseleccionar todas ellas es idéntico.

\subsubsection{Importación / Exportación}
Se provee funcionalidad de importar y exportar las anotaciones y mediciones de un modelo para que un profesor pueda compartirlas y para que los alumnos puedan enviar los ejercicios a un profesor y este las corrija.

Mediante el botón de exportar, descargaremos un fichero que contiene los elementos que tengamos en el momento de invocar dicha función. Nos aparecerá un menú en el que podremos elegir tanto la ubicación como el nombre del archivo.

Si lo que necesitamos es añadir a nuestro modelo un conjunto de elementos previamente creado, utilizaremos el botón de importar. Al pulsarlo, podremos ver un menú que nos permita seleccionar el archivo en cuestión.

\subsection{Subir modelos}
Finalmente, explicaremos cómo subir modelos. Sin embargo, esta vista puede que no esté permitida para algunos usuarios.

El uso de la misma es realmente sencilla; primero iremos a esta vista pulsando en <<Subir>> en la barra de navegación, después pulsaremos el botón de seleccionar archivo en la interfaz (típicamente se llamará <<Examinar>>), y nos aparecerá una ventana en la que seleccionar el modelo. Una vez seleccionado, pulsaremos el botón <<Subir>> y esperaremos a que el proceso termine. Una vez iniciada la subida aparecerá una barra que nos indicará el progreso, y cuando finalice el mismo se nos mostrará un mensaje avisándonos del éxito. Dicha pantalla puede ser vista en la figura~\ref{fig:viewer-upload-dragon}.
\imagenpersonalizada{viewer-upload-dragon}{Diálogo de subida de modelos}{1}



\bibliographystyle{plain}
\bibliography{bibliografiaAnexos}

\end{document}